\chapter[Lezione 21 --- 19 novembre 2026]{Lezione 21: continuità e classificazione delle discontinuità}
\label{chap:lezione-21}

\section{Continuità in un punto}
\label{sec:lez21-continuita}

La nozione di continuità collega il valore assunto da una
funzione in un punto al suo comportamento in prossimità
del punto stesso. Si distinguono i punti isolati e i punti
di accumulazione del dominio, introdotti nelle
\cref{def:lez08-isolato,def:lez07-accumulazione}.

\begin{definition}[Continuità in un punto]
\label{def:lez21-continuita}
Sia \(f\colon A\longrightarrow\R\), con \(A\subseteq\R\)
non vuoto, e sia \(x_0\in A\). La funzione \(f\) si dice
\emph{continua in \(x_0\)} se:
\begin{enumerate}
  \item \(x_0\) è isolato in \(A\); oppure
  \item \(x_0\) è di accumulazione per \(A\) e
  \[
    \lim_{x\to x_0}f(x)=f(x_0).
  \]
\end{enumerate}
La funzione è \emph{continua in \(B\subseteq A\)} se è
continua in ogni punto di \(B\).
\end{definition}

Nel primo caso la continuità è automatica: non è
necessario considerare il limite, che per la
\cref{def:lez08-limite} richiede un punto di accumulazione.
Nel secondo caso, la continuità equivale a
\[
  \lim_{x\to x_0}\bigl(f(x)-f(x_0)\bigr)=0.
\]

\section{Continuità delle funzioni elementari}
\label{sec:lez21-funzioni-elementari}

I limiti introdotti nelle lezioni precedenti permettono di
verificare la continuità delle principali funzioni.

\begin{proposition}[Continuità delle funzioni elementari]
\label{prop:lez21-continuita-elementari}
Valgono le seguenti proprietà:
\begin{enumerate}
  \item per ogni \(c\in\R\), la funzione costante
  \(f(x)=c\) è continua su \(\R\);
  \item per ogni \(n\in\N\), la funzione
  \(f(x)=x^n\) è continua su \(\R\);
  \item ogni funzione polinomiale
  \[
    P(x)=\sum_{j=0}^{m}a_jx^j,
    \qquad a_j\in\R,
  \]
  è continua su \(\R\);
  \item per ogni \(a>0\), la funzione esponenziale
  \(f(x)=a^x\) è continua su \(\R\);
  \item per ogni \(a>0\), con \(a\ne1\), la funzione
  logaritmica \(f(x)=\log_a x\) è continua su
  \((0,+\infty)\);
  \item le funzioni seno e coseno sono continue su
  \(\R\), mentre tangente e cotangente sono continue
  nei rispettivi domini:
  \[
  \begin{aligned}
    D(\tan)&=\R\setminus
      \left\{\frac\pi2+k\pi:k\in\Z\right\},\\
    D(\cot)&=\R\setminus
      \{k\pi:k\in\Z\}.
  \end{aligned}
  \]
\end{enumerate}
\end{proposition}

\clearpage

\begin{proof}[Dimostrazione del primo punto]
Per ogni \(x_0\in\R\), si ha
\[
  \lim_{x\to x_0}f(x)=c=f(x_0).
\]
Pertanto la funzione costante è continua su \(\R\).
\end{proof}

\begin{proof}[Dimostrazione del secondo punto]
Si procede per induzione su \(n\in\N\).

\emph{Passo base.}
Per \(n=0\), la funzione \(f(x)=x^0=1\)
è costante, quindi continua per il primo punto.

\emph{Passo induttivo.}
Sia \(n\ge0\) e si supponga che \(x\mapsto x^n\)
sia continua su \(\R\). Per ogni \(x_0\in\R\),
si sviluppi
\[
  x^{n+1}-x_0^{n+1}
  =x^n(x-x_0)+x_0(x^n-x_0^n).
\]
Per l'ipotesi induttiva,
\[
  \lim_{x\to x_0}x^n=x_0^n,
  \qquad
  \lim_{x\to x_0}(x-x_0)=0.
\]
Applicando i
\cref{thm:lez14-limite-somma,thm:lez14-limite-prodotto},
si ottiene
\[
\begin{aligned}
  \lim_{x\to x_0}(x^{n+1}-x_0^{n+1})
  &=x_0^n\cdot0+x_0\cdot0\\
  &=0.
\end{aligned}
\]
Pertanto
\[
  \lim_{x\to x_0}x^{n+1}=x_0^{n+1},
\]
quindi \(x\mapsto x^{n+1}\) è continua su \(\R\).
Per il \cref{thm:principio-induzione}, la proprietà
vale per ogni \(n\in\N\).
\end{proof}

\begin{proof}[Dimostrazione del terzo punto]
Sia \(x_0\in\R\). Per il secondo punto,
ogni funzione \(x\mapsto x^j\) è continua
in \(x_0\). Applicando i
\cref{thm:lez14-limite-somma,thm:lez14-limite-prodotto}
a una somma finita, si ottiene
\[
\begin{aligned}
  \lim_{x\to x_0}P(x)
  &=\sum_{j=0}^{m}a_j\lim_{x\to x_0}x^j\\
  &=\sum_{j=0}^{m}a_jx_0^j\\
  &=P(x_0).
\end{aligned}
\]
Pertanto ogni polinomio è continuo su \(\R\).
\end{proof}

\begin{proof}[Dimostrazione del quarto punto]
Per ogni \(a>0\) e \(x_0\in\R\), dal
\cref{thm:lez18-limite-esponenziale} segue
\[
  \lim_{x\to x_0}a^x=a^{x_0}.
\]
Pertanto la funzione esponenziale è continua
su \(\R\).
\end{proof}

\begin{proof}[Dimostrazione del quinto punto]
Per ogni \(a>0\), con \(a\ne1\), e \(x_0>0\),
dalla \cref{prop:lez20-limiti-logaritmo} segue
\[
  \lim_{x\to x_0}\log_a x=\log_a x_0.
\]
Pertanto la funzione logaritmica è continua
su \((0,+\infty)\).
\end{proof}

\begin{proof}[Dimostrazione del sesto punto]
Per ogni \(x_0\in\R\), dalla
\cref{prop:lez20-limiti-trigonometrici} segue
\[
  \lim_{x\to x_0}\sin x=\sin x_0,
  \qquad
  \lim_{x\to x_0}\cos x=\cos x_0.
\]
Pertanto seno e coseno sono continui su \(\R\).

Poiché
\[
  \tan x=\frac{\sin x}{\cos x},
  \qquad
  \cot x=\frac{\cos x}{\sin x},
\]
per il \cref{thm:lez14-limite-rapporto},
nei punti in cui i rispettivi denominatori
sono non nulli, si ha
\[
  \lim_{x\to x_0}\tan x=\tan x_0,
  \qquad
  \lim_{x\to x_0}\cot x=\cot x_0.
\]
Pertanto tangente e cotangente sono continue
nei rispettivi domini.
\end{proof}

\begin{example}[Una funzione non continua nell'origine]
\label{ex:lez21-seno-reciproco}
Si consideri
\[
  f(x):=
  \begin{cases}
    \sin(1/x), & x\ne0,\\
    0, & x=0.
  \end{cases}
\]
La funzione è definita su \(\R\) e \(f(0)=0\).
Per \(n\ge1\), si pongano
\[
  a_n:=\frac1{\frac\pi2+2\pi n},
  \qquad
  b_n:=\frac1{\frac{3\pi}2+2\pi n}.
\]
Entrambe le successioni tendono a \(0\), mentre
\[
  f(a_n)=1,\qquad f(b_n)=-1
  \qquad\forall n\ge1.
\]
Per il \cref{thm:lez17-teorema-ponte}, il limite
\(\lim_{x\to0}f(x)\) non esiste. Pertanto \(f\)
non è continua in \(0\), indipendentemente dal
valore assegnato in tale punto.
\end{example}

\section{Continuità unilaterale}
\label{sec:lez21-continuita-unilaterale}

La \cref{def:lez21-continuita} si può formulare anche
rispetto agli intorni unilateri della
\cref{def:lez14-limiti-unilateri}.

\begin{definition}[Continuità da destra e da sinistra]
\label{def:lez21-continuita-unilaterale}
Sia \(f\colon A\longrightarrow\R\) e sia \(x_0\in A\).
Se \(x_0\) è di accumulazione per \(A\) da destra,
\(f\) si dice \emph{continua da destra in \(x_0\)} se
\[
  \lim_{x\to x_0^+}f(x)=f(x_0).
\]
Se \(x_0\) è di accumulazione per \(A\) da sinistra,
\(f\) si dice \emph{continua da sinistra in \(x_0\)} se
\[
  \lim_{x\to x_0^-}f(x)=f(x_0).
\]
\end{definition}

Se \(x_0\) è di accumulazione da entrambi i lati,
per il \cref{thm:lez14-limite-unilateri} la continuità
in \(x_0\) equivale alla continuità da destra e da
sinistra nello stesso punto.

\clearpage

\begin{example}[Continuità della funzione parte intera]
\label{ex:lez21-continuita-floor}
Si consideri la funzione
\[
  f\colon\R\longrightarrow\R,
  \qquad f(x):=\lfloor x\rfloor.
\]
Per la \cref{def:floor}, se \(m\in\Z\), allora
\[
  f(m)=m,
  \qquad
  \lim_{x\to m^+}f(x)=m,
  \qquad
  \lim_{x\to m^-}f(x)=m-1.
\]
Pertanto \(f\) è continua da destra, ma non da
sinistra, in ogni \(m\in\Z\).

Se \(x_0\in\R\setminus\Z\), posto
\(m:=\lfloor x_0\rfloor\), si ha
\(x_0\in(m,m+1)\). In un intorno di \(x_0\)
la funzione vale costantemente \(m\), quindi
\[
  \lim_{x\to x_0}f(x)=m=f(x_0).
\]
La funzione parte intera è dunque continua in
\(\R\setminus\Z\) e discontinua in ogni intero,
come illustrato nella \cref{fig:floor}.
\end{example}

\section{Classificazione delle discontinuità}
\label{sec:lez21-discontinuita}

Sia \(f\colon A\longrightarrow\R\), con
\(A\subseteq\R\) non vuoto, e sia \(x_0\in A\)
un punto di accumulazione per \(A\). Se \(f\)
non è continua in \(x_0\), si distinguono
i seguenti casi.

\begin{definition}[Tipi di discontinuità]
\label{def:lez21-tipi-discontinuita}
\begin{enumerate}
  \item La discontinuità è \emph{eliminabile} se
  esiste \(\ell\in\R\) tale che
  \[
    \lim_{x\to x_0}f(x)=\ell\ne f(x_0).
  \]
  \item La discontinuità è \emph{di prima specie}
  se esistono e sono finiti entrambi i limiti unilateri,
  ma sono distinti:
  \[
    \lim_{x\to x_0^-}f(x)=\ell_-,\qquad
    \lim_{x\to x_0^+}f(x)=\ell_+,
    \qquad\ell_-\ne\ell_+.
  \]
  \item La discontinuità è \emph{di seconda specie}
  negli altri casi in cui \(f\) non ammette limite
  finito in \(x_0\); in particolare, quando
  almeno uno dei limiti unilateri è infinito
  o non esiste.
\end{enumerate}
\end{definition}

Nei punti di prima specie si suppone che il dominio
ammetta punti arbitrariamente vicini a \(x_0\) da
entrambi i lati, affinché i due limiti unilateri
siano definiti. Gli esempi schematici della
\cref{fig:lez21-tipi-discontinuita} illustrano le
tre situazioni.

\begin{figure}[htbp]
  \centering
  
\begin{tikzpicture}[unipd/diagram,x=1cm,y=1cm]
  \def\xmin{-1.6}
  \def\xmax{1.6}
  \def\ymin{-0.2}
  \def\ymax{2.7}

  \begin{scope}[shift={(0,0)}]
    \draw[->,unipd-muted] (\xmin,0)--(\xmax,0)
      node[right] {$x$};
    \draw[->,unipd-muted] (0,\ymin)--(0,\ymax)
      node[above] {$y$};

    \draw[very thick,unipd-primary,domain=-1.35:1.35,samples=100]
      plot (\x,{1.2+0.52*\x*\x});

    \fill[white] (0,1.2) circle (2.5pt);
    \draw[unipd-primary,line width=.8pt]
      (0,1.2) circle (2.5pt);
    \fill[unipd-primary] (0,2.1) circle (2.5pt);

    \node[below] at (0,\ymin) {$x_0$};
    \node[align=center,font=\sffamily\footnotesize]
      at (0,-1.05) {Discontinuità\\eliminabile};
  \end{scope}

  \begin{scope}[shift={(4.8,0)}]
    \draw[->,unipd-muted] (\xmin,0)--(\xmax,0)
      node[right] {$x$};
    \draw[->,unipd-muted] (0,\ymin)--(0,\ymax)
      node[above] {$y$};

    \draw[very thick,unipd-primary,domain=-1.35:0,samples=60]
      plot (\x,{1+0.33*\x});
    \draw[very thick,unipd-primary,domain=0:1.35,samples=60]
      plot (\x,{1.9+0.33*\x});

    \fill[white] (0,1) circle (2.5pt);
    \draw[unipd-primary,line width=.8pt]
      (0,1) circle (2.5pt);
    \fill[white] (0,1.9) circle (2.5pt);
    \draw[unipd-primary,line width=.8pt]
      (0,1.9) circle (2.5pt);
    \fill[unipd-primary] (0,1.45) circle (2.5pt);

    \node[below] at (0,\ymin) {$x_0$};
    \node[align=center,font=\sffamily\footnotesize]
      at (0,-1.05) {Discontinuità\\di prima specie};
  \end{scope}

  \begin{scope}[shift={(9.6,0)}]
    \draw[->,unipd-muted] (\xmin,0)--(\xmax,0)
      node[right] {$x$};
    \draw[->,unipd-muted] (0,\ymin)--(0,\ymax)
      node[above] {$y$};

    \draw[densely dashed,unipd-muted]
      (0,.75)--(0,2.6);

    \draw[very thick,unipd-primary,domain=-1.35:-.15,samples=100]
      plot (\x,{.4-.3/\x});
    \draw[very thick,unipd-primary,domain=.15:1.35,samples=100]
      plot (\x,{.4+.3/\x});

    \fill[unipd-primary] (0,.55) circle (2.5pt);

    \node[below] at (0,\ymin) {$x_0$};
    \node[align=center,font=\sffamily\footnotesize]
      at (0,-1.05) {Discontinuità\\di seconda specie};
  \end{scope}
\end{tikzpicture}

  \caption{Discontinuità eliminabile, di prima specie e
  di seconda specie. I cerchi vuoti indicano i valori
  limite non assunti nel punto, mentre i punti pieni
  indicano il valore della funzione.}
  \label{fig:lez21-tipi-discontinuita}
\end{figure}

\begin{example}[Discontinuità eliminabile e non eliminabile]
\label{ex:lez21-prolungamenti}
La funzione
\[
  h\colon\R\setminus\{0\}\longrightarrow\R,
  \qquad h(x):=\frac{\sin x}{x},
\]
non è definita in \(0\), ma per il
\cref{thm:lez20-limite-seno} si ha
\[
  \lim_{x\to0}h(x)=1.
\]
Assegnando il valore \(h(0):=1\), si ottiene un
prolungamento continuo su \(\R\). Assegnando
invece un valore diverso da \(1\), si ottiene una
discontinuità eliminabile.

Al contrario, per \(g(x):=1/|x|\), definita per
\(x\ne0\), si ha
\[
  \lim_{x\to0}g(x)=+\infty.
\]
Nessun valore reale assegnato a \(g(0)\) rende
la funzione continua nell'origine: ogni tale
prolungamento presenta una discontinuità
di seconda specie.
\end{example}

\section{Operazioni con funzioni continue}
\label{sec:lez21-operazioni-continuita}

Le operazioni tra funzioni continue si studiano
mediante i teoremi sulle operazioni dei limiti
introdotti nella \cref{sec:lez14-operazioni-limiti}.

\begin{proposition}[Continuità di somma, prodotto e rapporto]
\label{prop:lez21-operazioni-continuita}
Siano \(f,g\colon A\longrightarrow\R\), con
\(A\subseteq\R\) non vuoto, e sia \(x_0\in A\).
Se \(f\) e \(g\) sono continue in \(x_0\), allora
anche \(f+g\) e \(fg\) sono continue in \(x_0\).
Se inoltre \(g(x_0)\ne0\), il rapporto \(f/g\)
è continuo in \(x_0\), nel proprio dominio.
\end{proposition}

\begin{proof}
Se \(x_0\) è isolato in \(A\), la continuità
segue dalla \cref{def:lez21-continuita}.

Sia dunque \(x_0\) di accumulazione per \(A\).
Per ipotesi,
\[
  \lim_{x\to x_0}f(x)=f(x_0),\qquad
  \lim_{x\to x_0}g(x)=g(x_0).
\]
Per i
\cref{thm:lez14-limite-somma,thm:lez14-limite-prodotto},
\[
  \lim_{x\to x_0}(f+g)(x)=f(x_0)+g(x_0),
\]
\[
  \lim_{x\to x_0}(fg)(x)=f(x_0)g(x_0).
\]
Se \(g(x_0)\ne0\), per la continuità di \(g\)
esiste un intorno \(U_{x_0}\) tale che
\(g(x)\ne0\) per ogni \(x\in U_{x_0}\cap A\).
Pertanto, per il \cref{thm:lez14-limite-rapporto},
si ha inoltre
\[
  \lim_{x\to x_0}\frac{f(x)}{g(x)}
  =\frac{f(x_0)}{g(x_0)}.
\]
Le tre uguaglianze verificano la definizione
di continuità.
\end{proof}

\section{Continuità della funzione composta}
\label{sec:lez21-continuita-composta}

\begin{proposition}[Continuità della composizione]
\label{prop:lez21-continuita-composta}
Siano \(f\colon A\longrightarrow\R\) e
\(g\colon B\longrightarrow A\), con \(A,B\subseteq\R\)
non vuoti. Sia \(t_0\in B\) di accumulazione
per \(B\) e si ponga \(x_0:=g(t_0)\in A\).
Se \(x_0\) è di accumulazione per \(A\),
\(g\) è continua in \(t_0\) e \(f\) è continua
in \(x_0\), allora \(f\circ g\) è continua
in \(t_0\).
\end{proposition}

\begin{proof}
Per la continuità delle due funzioni,
\[
  \lim_{t\to t_0}g(t)=g(t_0)=x_0,
  \qquad
  \lim_{x\to x_0}f(x)=f(x_0).
\]
Poiché \(x_0\in A\) e \(f(x_0)\) coincide con
il limite di \(f\), è soddisfatta la seconda
condizione del \cref{thm:lez16-limite-sostituzione}.
Pertanto
\[
  \lim_{t\to t_0}f(g(t))
  =f(x_0)=f(g(t_0)),
\]
e la funzione composta è continua in \(t_0\).
\end{proof}

\begin{example}[Continuità di una funzione composta]
\label{ex:lez21-seno-esponenziale}
La funzione \(x\mapsto e^x\) è continua su \(\R\)
per il \cref{thm:lez18-limite-esponenziale},
mentre \(x\mapsto\sin x\) è continua su \(\R\)
per la \cref{prop:lez20-limiti-trigonometrici}.
Dalla \cref{prop:lez21-continuita-composta} segue
che \(x\mapsto\sin(e^x)\) è continua su \(\R\).
\end{example}

\begin{exercise}[Continuità di una funzione definita a tratti]
\label{ex:lez21-continuita-tratti}
Determinare se esiste \(a\in\R\) tale che la funzione
\[
  f(x):=
  \begin{cases}
    \dfrac{\sin x}{x}, & x>0,\\[4pt]
    a, & x=0,\\[4pt]
    \log(1+|x|), & x<0
  \end{cases}
\]
sia continua su \(\R\).
\end{exercise}

\begin{proof}[Soluzione dell'esercizio]
Per \(x>0\), la funzione è continua per la
\cref{prop:lez21-operazioni-continuita}, essendo
\(x\ne0\). Per \(x<0\), si ha \(|x|=-x\),
e quindi la funzione è continua come composizione
di \(x\mapsto1-x\) e \(u\mapsto\log u\).

Resta da esaminare \(x_0=0\). Dal
\cref{thm:lez20-limite-seno} e dalla
\cref{prop:lez20-limiti-logaritmo} segue
\[
  \lim_{x\to0^+}f(x)=1,
  \qquad
  \lim_{x\to0^-}f(x)=0.
\]
Poiché i limiti unilateri sono distinti,
per il \cref{thm:lez14-limite-unilateri}
non esiste il limite in \(0\). Pertanto
\emph{non esiste alcun \(a\in\R\)} che renda
\(f\) continua su \(\R\): in \(0\) la funzione
presenta una discontinuità di prima specie.
\end{proof}

\section{Continuità dell'applicazione inversa}
\label{sec:lez21-continuita-inversa}

\begin{theorem}[Continuità dell'applicazione inversa]
\label{thm:lez21-continuita-inversa}
Sia \(I\subseteq\R\) un intervallo e sia
\(f\colon I\longrightarrow\R\) strettamente monotona.
Allora l'applicazione inversa
\[
  f^{-1}\colon\operatorname{Im}(f)\longrightarrow I
\]
è continua su \(\operatorname{Im}(f)\).
\end{theorem}

La stretta monotonia garantisce l'iniettività
(\cref{prop:lez10-monotonia-inversa}) e quindi
l'esistenza dell'inversa della \cref{def:lez09-inversa}.
Il teorema non richiede che \(f\) sia continua.
La dimostrazione non viene sviluppata in questa sede.
